% !PRINTED none
% Unnumbered section-opening page (top portion of the printed page is blank, no running head/page number visible); begins the errata section with the centered heading "Berichtigungen zum ersten Bande.", its errata table, then the heading "Berichtigungen zum zweiten Bande." with its errata table continuing into an unnumbered displayed formula.

\begin{center}
\textbf{Berichtigungen zum ersten Bande.}
\end{center}

\noindent (Kleinere Versehen, die der Leser leicht selbst verbessert, sind hier nicht aufgeführt.)

\begin{center}
\rule{3cm}{0.4pt}
\end{center}

\bigskip
\noindent Seite 185, Zeile 4 v. o. lies $(m-s)(n-s)+s$ statt $(m-s)(n-s)$ nach
einigen leicht ersichtlichen Änderungen in den letzten Zeilen von
Seite 184.\\
\glqq{} 186, Formel (4) lies $n\mu + m\nu$ statt $n\nu + m\mu$.\\
\glqq{} 195, Zeile 7 v. o. lies zwei statt drei.\\
\glqq{} 198, \glqq{} 13 v. o. \glqq{} $U_1^2$ statt $U_1$.\\
\glqq{} 215, \glqq{} 7 v. u. \glqq{} $\lambda : r^2$ statt $\lambda r^2$.\\
\glqq{} 221, Formel (13) und (14) lies $r^{-1}$ statt $r$.\\
\glqq{} 246, Zeile 6 v. u. ist $O_1$ zu streichen.\\
\glqq{} 320, Formel 7 das zweite $\Phi(z)\theta(z)^2$ durch $\Phi(x)\theta(x)^2$ zu ersetzen.\\
\glqq{} 339, Zeile 4 v. u. lies reellem statt imaginärem.\\
\glqq{} 445, \glqq{} 19 v. o. \glqq{} unlösbar statt lösbar.\\
\glqq{} 487, \glqq{} 13 v. o. \glqq{} $(\mathit{mod}\,p)$ statt $(\mathit{mod}\,m)$.\\
\glqq{} 541, \glqq{} 9 v. o. \glqq{} $\mu$ statt $m$.\\
\glqq{} 552, \glqq{} 11 v. o. \glqq{} $\psi'/\pi$ statt $\psi/\pi$.\\
\glqq{} 657, \glqq{} 7 v. u. ist hinter \glqq Normalgleichung\grqq{} einzuschalten: mit
reellen Wurzeln.

\begin{center}
\rule{3cm}{0.4pt}
\end{center}

\begin{center}
\textbf{Berichtigungen zum zweiten Bande.}
\end{center}

\begin{center}
\rule{3cm}{0.4pt}
\end{center}

\bigskip
\noindent Seite 231, 232. Bei der Anwendung, die hier von dem Satze 4 gemacht ist,
wurde übersehen, daß die Invarianten selbst zum Rationalitätsbereich
 gehören, daß also die Darstellung von $J$ als rationale
Funktion von $\theta$ eine reine Identität ($J = J$) ist. Demnach ist
der Absatz Seite 231 unten von \glqq Aus diesem Satze\grqq{} bis 232 oben
\glqq unentschieden bleiben\grqq{} zu streichen.\\
\glqq{} 640, Zeile 14 v. o. lies $F'(\tau)$ statt $F'(\tau')$.\\
\glqq{} 697, \glqq{} 13 v. u.: Die hier folgende Schlußweise muß folgendermaßen
abgeändert werden:

\bigskip
\noindent Aus (13) ergibt sich durch Multiplikation mit $g_s$ und Summation:

\[
-n\varkappa \Sigma \delta_s g_s + \alpha u < \Sigma g_s \lambda_s < \Sigma \delta_s g_s + \alpha u.
\]
